% Source: 01 Properties_of_Matter 26(1).docx; formatting in layout.tex.
\ResponseGap

\HOne{The Chemical Reaction}

Chemistry studies substances, their properties, and the changes they undergo. One particularly important aspect of chemistry is the formation of new substances. Such a process is called a \textbf{chemical reaction}. In a chemical reaction, one or more starting substances form \textbf{new substances} with \textbf{different characteristic properties}. The starting substances are called reactants, and the substances formed are called products.


Examples of chemical reactions include the combustion of paper, wood, petrol, or natural gas, the rusting of iron, the changing colour of leaves in autumn, the digestion of food, and the discharge of a battery. In each case, substances with new properties are formed.


If \textbf{no new substance} is formed, the change is called a \textbf{physical }\textbf{process} rather than a chemical reaction. During a physical process, a substance may change its state, shape, or appearance, but it remains the same substance. For example, when water freezes, it becomes solid, but it is still water.


Another characteristic of a chemical reaction is the \textbf{transfer of energy} between the reacting substances and their surroundings. Based on the overall direction of this energy transfer, chemical reactions can be classified as \textbf{exothermic} or \textbf{endothermic}.


Many chemical reactions do not start on their own, even when the reactants are in contact. A match must be struck, the fuel--air mixture in a petrol engine must be ignited by a spark before combustion can begin.


The initial amount of energy required to start a chemical reaction is called the \textbf{activation energy}.

\begin{InfoBox}
    \BoxHeading{Exothermic reaction}
    The substances which react give off energy to their surroundings. \newline \textit{Examples: combustion, explosions, discharging of a battery}
\end{InfoBox}

\ResponseGap[200pt]

\newpage

\begin{InfoBox}
\BoxHeading{Endothermic reaction}
The substances which react have to \textbf{take up energy} from their surroundings.\newline \textit{Examples: making toast, photosynthesis in green plants, charging of a battery}
\end{InfoBox}

\ResponseGap[200pt]

\HFour{Examples of Chemical Reaction Equations}

\ResponseGap[90pt]

\HFour{Chemical Reactions on a Particle Level}

\ResponseGap[90pt]

Based on the particle model, we have to assume that material particles can be formed into different types of particles. In a uniform reaction, the same types of reactant particles form the same types of product particles.


For a reaction between substances A and B to take place, their particles must come into contact. In solids, most particles are fixed in place, so the reaction usually occurs mainly where the surfaces of the two substances touch.


\begin{center}% Replaces images/image17.jpeg (original retained)
\resizebox{0.75\linewidth}{!}{\SolidSurfaceReactionFigure}\end{center}


In liquids and gases, particles can move and mix more freely. This often makes a reaction faster. Dissolving a solid can have a similar effect because its particles can come into contact more easily with the particles of another reactant.

